\subsection{Линейный график технологического процесса}

В целях планирования сроков выполнения работ и определения их взаимосвязи используется метод сетевого планирования. Данный метод позволяет представить процесс разработки программного средства в виде ориентированного графа, отражающего последовательность выполнения работ и зависимости между ними.

В сетевой модели каждая работа характеризуется кодом, состоящим из номера начального и конечного событий. События отражают моменты завершения отдельных этапов разработки, а работы — процессы, необходимые для перехода от одного события к другому. Такая форма представления позволяет определить критический путь проекта, оценить общую продолжительность разработки, а также выявить возможные резервы времени.

На основе ранее определённых работ и их трудоёмкости были сформированы коды работ и рассчитаны временные параметры сетевого графа. Для каждой работы определены её продолжительность, ранние и поздние сроки наступления событий, а также временные резервы. Результаты расчётов параметров сетевого графа приведены в таблице~\ref{tab:work_codes}.


{
\fontsize{10}{10}\selectfont
\setlength{\tabcolsep}{3pt}

\begin{longtable}{|c|p{3cm}|p{3cm}|p{1.5cm}|p{1.5cm}|p{1.2cm}|p{1.2cm}|p{2.9cm}|}
\caption{Соответствие событий и работ}\label{tab:work_codes} \\
\hline
№ & Событие & Работа & Т (ч) & Т (дн) & Код  & Тип & Исполнитель \\ \hline
\endfirsthead

\caption*{Продолжение таблицы \ref{tab:work_codes}}\\ \hline
№ & Событие & Работа & Т (ч) & Т (дн) & Код  & Тип & Исполнитель \\ \hline
\endhead

\hline
\endfoot

\hline
\endlastfoot

1 & Начало работ & Анализ предметной области и существующих решений & 28 & 3.5 & 1-2 & ОН & Аналитик \\ \hline

2 & Завершение анализа предметной области & Формирование требований к системе & 18 & 2.3 & 2-3 & ОН & Аналитик \\ \hline

\multirow{2}{*}{3} & \multirow{2}{*}{\parbox{3cm}{Утверждение технического задания}}
& Проектирование общей архитектуры системы & 40 & 5 & 3-4 & ОН & Инженер \\ \cline{3-8}
& & Подготовка и предобработка данных & 56.0  & 7.0 & 3-6 & ОН & Аналитик \\ \hline

4 & Проектирование общей архитектуры системы завершено & Проектирование структуры данных & 19.6  & 2.5 & 4-5 & ОН & Инженер \\ \hline

5 & Проектирование структуры данных завершено & Реализация модуля сбора и обработки логов & 112 & 14.0 & 5-9 & ОН & Инженер \\ \hline

6 & Подготовка и предобработка данных завершена & Разработка моделей НС для обнаружения аномалий & 67.2 & 8.4 & 6-7 & ОН & Аналитик \\ \hline

7 & Разработка моделей НС завершена & Обучение моделей НС & 50.4 & 6.3 & 7-8 & ОН & Аналитик \\ \hline

8 & Обучение моделей НС завершено & Оценка и сравнение моделей & 28  & 3.5 & 8-10 & ОН & Аналитик \\ \hline

9 & Реализация модуля сбора и обработки логов завершена & Реализация модуля анализа и обнаружения аномалий & 64.0  & 8.0 & 9-10 & ОН & Инженер \\ \hline

10 & Реализация модуля анализа и обнаружения аномалий завершена & Интеграция модели НС в систему & 44.0  & 5.5  & 10-11 & ОН & Инженер \\ \hline

11 & Интеграция модели НС в систему завершена & Настройка взаимодействия компонентов системы & 44.0  & 5.5  & 11-12 & ОН & Инженер \\ \hline

12 & Настройка взаимодействия компонентов системы завершена & Функциональное тестирование & 176  & 22  & 12-13 & ОН & Инженер \\ \hline

13 & Функциональное тестирование завершено & Исправление ошибок и оптимизация & 112.0 & 14.0  & 13-14 & ОН & Инженер \\ \hline

14 & Исправление ошибок и оптимизация завершены & Подготовка документации & 168  & 21 & 14-15 & ОН & Инженер \\ \hline

15 & Завершение проекта & - & -  & - & - & - & - \\ \hline

\end{longtable}
}